% ---
% id: meat-core/frank-wolfe/prop:relaxed-set-properties
% kind: prop
% title: "Properties of the relaxed feasible set"
% summary: >-
%   The relaxed feasible sets grow with the window, are convex and compact,
%   and keep every hypothesis of this module that names the feasible set,
%   positivity carrying them; the objective falls below zero there by at most
%   an explicit floor. Reading the set at a window can only lower the value.
% prominence: standard
% status: implemented
% uses: [meat-core/thm:branch-formula, meat-core/lem:branch-maps, foundations/lem:pinchmono, def:relaxed-feasible-set, meat-core/lem:cone-concavity, meat-core/thm:convexity, meat-core/thm:reduction, meat-core/cor:reduction-constrained, meat-core/thm:choi, meat-core/prop:value-range, meat-core/def:support-frame, meat-core/lem:choi-factor, meat-core/def:feasible, meat-core/lem:joint-sectors, meat-core/def:admissible-declaration]
% mentions: [meat-core/def:untrusted-choi, meat-core/cor:concavity, meat-core/def:maps, meat-core/def:declaration, thm:firstorder, foundations/req:norm-conversion, foundations/req:tolerance-scope, req:iterate-checks, req:iterate-repair]
% assumes: []
% refs: []
% provenance: newly-derived
% implements: certiqs_key.solvers.frank_wolfe.joint_driver:_check_point
% lean:
% ---

The relaxed feasible sets keep every hypothesis of this module that names
the feasible set, and positivity of the variable is what carries them:
(a) and (b) are the monotonicity, convexity and compactness of the
sets, (c) the facts the objective needs on them, (c$'$) the implication
behind each normalisation clause, (c$''$) what relaxing the normalisation
costs, (d), (d$'$) that reading the set at a tolerance is conservative,
and (e) that all of this holds in the general format.
The window at which a run checks the positivity of a point it consumes is
a number this module declares and records
(\xref{foundations/req:tolerance-scope}\,(ii)); item (c) \emph{Fourth}
gives the reason it reads $0$ (Requirements~\ref{req:iterate-checks}
and~\ref{req:iterate-repair}).

\begin{proposition}\label{prop:relaxed-set-properties}
Let $\Decl$ satisfy Definition~\ref{def:admissible-declaration}, let $\varepsilon\ge0$, and
let $\Feas(\varepsilon)$ and $\Feas^{J}(\varepsilon)$ be as in
Definition~\ref{def:relaxed-feasible-set}.
\begin{enumerate}[label=(\alph*),leftmargin=2.4em]
\item \emph{(Monotone.)} For $0\le\varepsilon\le\varepsilon'$,
  $\Feas(\varepsilon)\subseteq\Feas(\varepsilon')$ and
  $\Feas^{J}(\varepsilon)\subseteq\Feas^{J}(\varepsilon')$; in particular
  $\Feas\subseteq\Feas(\varepsilon)$ and $\Feas^{J}\subseteq\Feas^{J}(\varepsilon)$
  for every $\varepsilon\ge0$.
\item \emph{(Convex and compact.)} $\Feas(\varepsilon)$ and
  $\Feas^{J}(\varepsilon)$ are convex and compact. Every
  $\omega\in\Feas(\varepsilon)$ has $\tr\omega\le1+d_A\varepsilon$ in
  trusted-source mode and $\tr\omega\le1+\varepsilon$ in untrusted-source
  mode, and every $J\in\Feas^{J}(\varepsilon)$ has
  $\sum_s\tr J_s\le d_{A_0}(1+\varepsilon)$, a bound that reads the coupling
  clause and never the normalisation clause.
\item \emph{(What positivity carries.)} $\Feas(\varepsilon)\subseteq\Pos{A\Qsys}$
  and $\chi_{\Qsys}\big(\Feas^{J}(\varepsilon)\big)\subseteq\Pos{A\Qsys}$.
  \emph{First}, the bounds $0\le K_{\hat c}(\omega)\le P_{\hat c}(\omega)$
  hold on the whole cone. \emph{Second}, $K_{\hat c}$ is concave on the
  sector-block-diagonal part of $\Pos{A\Qsys}$. \emph{Third},
  Theorem~\ref{thm:convexity} holds verbatim with $\Feas(\varepsilon)$ in
  place of $\Feas$, and $O^{J}_{\alpha}$ is convex and lower semicontinuous
  on $\acSet\times\Feas^{J}(\varepsilon)$. \emph{Fourth}, none of this
  survives a relaxation of positivity: the objective factors through
  $X\mapsto X^{p}$, which is undefined off the cone. One fact of the exact
  set does not survive: for $\varepsilon>0$ and $\omega\in\Feas(\varepsilon)$,
  $\sum_{\hat c}P_{\hat c}(\omega)=\tr\omega$ for an exactly admissible
  $\{F_v\}$, which the relaxed
  normalisation does not pin to $1$, so $\Phimap[\omega]$ need not be a
  probability vector and the range fact $\hlow\ge0$ of
  Proposition~\ref{prop:value-range}\,(i) is not available on
  $\Feas(\varepsilon)$ or on $\Feas^{J}(\varepsilon)$; it holds on the exact
  sets $\Feas$ and $\Feas^{J}$.
\item[(c$'$)] \emph{(Each normalisation clause is implied by the clause
  before it.)} Let $J=(J_s)_s$ be per-block positive semidefinite with
  $\|R\|_{\max}\le\varepsilon$, where
  $R\coloneqq\sum_s\tr_{\Qsys_s}J_s-\idon{A_0}$, and put
  $G\coloneqq\iota^{\dagger}\tau_A\iota=W_{\tau}^{\dagger}W_{\tau}$. Then
  $G$ is a state,
  \begin{equation}\label{eq:choi-trace-residual}
    \tr[\chi_{\Qsys}(J)]
    =\tr\Big[W_{\tau}^{\dagger}W_{\tau}\,
      \textstyle\sum_s\tr_{\Qsys_s}J_s\Big]
    =\tr\big[G(\idon{A_0}+R)\big]
    =1+\tr[GR],
  \end{equation}
  and $\big|\tr[\chi_{\Qsys}(J)]-1\big|\le d_{A_0}\varepsilon$; in
  particular $\tr[\chi_{\Qsys}(J)]=1$ holds at a point of nonzero residual
  $R$ exactly where $\tr[GR]=0$. In the state form in trusted-source mode,
  $\|\omega_A-\tau_A\|_{\max}\le\varepsilon$ implies
  $|\tr\omega-1|\le d_A\varepsilon$. In untrusted-source mode the
  normalisation clause is the relaxed coupling condition itself.
\item[(c$''$)] \emph{(What relaxing the normalisation costs.)} At a point
  of $\Feas^{J}(\varepsilon)$ put
  $\varsigma\coloneqq\tr[\chi_{\Qsys}(J)]-1$, so that
  $|\varsigma|\le d_{A_0}\varepsilon$ and
  $\sum_{\hat c}P^{J}_{\hat c}(J)=1+\varsigma$, and, where $1+\varsigma>0$,
  $\bar{\mathbf P}\coloneqq\mathbf P/(1+\varsigma)$, a probability vector.
  Then for every probability vector $\mathbf q$
  \begin{equation}\label{eq:penaltyshift-fw}
    \sum_{\hat c}q(\hat c)\log_2\frac{q(\hat c)}{P^{J}_{\hat c}}
    =\sum_{\hat c}q(\hat c)\log_2\frac{q(\hat c)}{(1+\varsigma)\bar P_{\hat c}}
    =\sfD(\mathbf q\|\bar{\mathbf P})-\log_2(1+\varsigma),
  \end{equation}
  and, provided $d_{A_0}\varepsilon<1$,
  \begin{equation}\label{eq:relaxed-objective-floor}
    O^{J}_{\alpha}(\mathbf q,J)\;\ge\;
    -\frac{\alpha}{\alpha-1}\,\log_2\!\big(1+d_{A_0}\varepsilon\big)
    \qquad\text{on }\acSet\times\Feas^{J}(\varepsilon),
  \end{equation}
  whose right-hand side is $0$ at $\varepsilon=0$ and negative otherwise.
  The state form carries the same statement with $d_A\varepsilon$ in
  trusted-source mode and $\varepsilon$ in untrusted-source mode in place of
  $d_{A_0}\varepsilon$. At $d_{A_0}\varepsilon\ge1$ the identity
  \eqref{eq:penaltyshift-fw} still holds wherever $1+\varsigma>0$ but
  carries no bound.
\item[(d)] \emph{(Reading the set at a tolerance is conservative.)} With
  $\hlow[\Feas(\varepsilon)]\coloneqq\inf_{\mathbf q\in\acSet}
  \inf_{\omega\in\Feas(\varepsilon)}O_{\alpha}(\mathbf q,\omega)$,
  \begin{equation}\label{eq:feas-tol-mono}
    \varepsilon\le\varepsilon'\ \Longrightarrow\
    \hlow\big[\Feas(\varepsilon')\big]\;\le\;\hlow\big[\Feas(\varepsilon)\big]
    \;\le\;\hlow\big[\Feas(0)\big]=\hlow\;\le\;\hqkd .
  \end{equation}
\item[(d$'$)] \emph{(The Choi-form quantity.)} With
  \begin{equation}\label{eq:hlow-feasJ}
    \hlow\big[\Feas^{J}(\varepsilon)\big]\coloneqq
    \inf_{\mathbf q\in\acSet}\ \inf_{J\in\Feas^{J}(\varepsilon)}
    O^{J}_{\alpha}(\mathbf q,J),
  \end{equation}
  the two relaxed quantities agree at $\varepsilon=0$, and
  \begin{equation}\label{eq:feas-tol-mono-J}
    \varepsilon\le\varepsilon'\ \Longrightarrow\
    \hlow\big[\Feas^{J}(\varepsilon')\big]\;\le\;
    \hlow\big[\Feas^{J}(\varepsilon)\big]\;\le\;
    \hlow\big[\Feas^{J}(0)\big]\;=\;\hlow\;\le\;\hqkd .
  \end{equation}
  At $\varepsilon>0$ no relation between $\hlow[\Feas^{J}(\varepsilon)]$
  and $\hlow[\Feas(\varepsilon)]$ is asserted.
\item[(e)] \emph{(The general format.)} For a declaration in the general
  format, with both sets read as Definition~\ref{def:relaxed-feasible-set}
  reads them there and $\chi_{\Qsys}$ the assignment
  $(J_{(\mathsf t,s)})\mapsto(\omega^{(\mathsf t,s)})$ of
  Theorem~\ref{thm:choi}, (a) to (d$'$) hold as stated. In (c$'$) the
  residual is $R=\bigoplus_{\mathsf t}R_{\mathsf t}$, $G=\bigoplus_{\mathsf
  t}G_{\mathsf t}$ with $G_{\mathsf t}\coloneqq W_{\mathsf t}^{\dagger}
  W_{\mathsf t}$, and
  $\tr[\chi_{\Qsys}(J)]=1+\sum_{\mathsf t}\tr[G_{\mathsf t}R_{\mathsf t}]$.
\end{enumerate}
\end{proposition}

\begin{proof}
\emph{Idea.} Every relaxed condition weakens as the window grows, which
gives (a), and is an affine preimage of a closed convex set, which with
the trace bounds gives (b). The facts of (c) read only positivity
of the variable, which the relaxation keeps, and transfer to the Choi form
along the linear map $\chi_{\Qsys}$. Item (c$'$) is a trace estimate on the
coupling residual, (c$''$) a rescaling of the statistics vector, and (d),
(d$'$) compare infima over nested sets.

\emph{(a).} For $0\le\varepsilon\le\varepsilon'$ each condition of
\eqref{eq:feas-eps} at $\varepsilon$ implies the same condition at
$\varepsilon'$, and so does each condition of the untrusted-source and the
Choi forms, so $\Feas(\varepsilon)\subseteq\Feas(\varepsilon')$ and
$\Feas^{J}(\varepsilon)\subseteq\Feas^{J}(\varepsilon')$. At
$\varepsilon=0$ the sets are $\Feas$ and $\Feas^{J}$
(Definition~\ref{def:relaxed-feasible-set}), which gives the inclusions
for every $\varepsilon\ge0$.

\emph{(b).} Each relaxed condition is convex:
$\lmin(B-A(\omega))\ge-\varepsilon$ reads
$B-A(\omega)+\varepsilon\idon{}\succeq0$, an affine preimage of the
positive semidefinite cone; $\|\cdot\|_{\max}\le\varepsilon$ is an affine
preimage of a norm ball; each normalisation clause is the preimage of a
closed interval under the affine map $\omega\mapsto\tr\omega$,
respectively $J\mapsto\tr\chi_{\Qsys}(J)$, $\chi_{\Qsys}$ being linear;
and $\omega\succeq0$ and the block structure are as for $\Feas$. Each is
closed, so $\Feas(\varepsilon)$ is closed, and it is bounded because
$\omega\succeq0$ with $\tr\omega\le1+d_A\varepsilon$ in trusted-source
mode and $\tr\omega\le1+\varepsilon$ in untrusted-source mode: the relaxed
normalisation enters here as a bound on the trace in place of its exact
value $1$. $\Feas^{J}(\varepsilon)$ is closed for the same reasons and
bounded because $J_s\succeq0$ with
$\sum_s\tr J_s=\tr\big[\sum_s\tr_{\Qsys_s}J_s\big]
\le d_{A_0}\,(1+\varepsilon)$: the Choi input register $A_0$ has dimension
$d_{A_0}$ (Definition~\ref{def:support-frame}), and
$|\tr X|\le d\,\|X\|_{\max}$ for a $d\times d$ matrix, applied to
$X=\sum_s\tr_{\Qsys_s}J_s-\idon{A_0}$, gives
$|\tr[\sum_s\tr_{\Qsys_s}J_s]-d_{A_0}|\le d_{A_0}\,\varepsilon$. That
bound reads the coupling clause and never the normalisation clause.

\emph{(c), state form.} $\Feas(\varepsilon)\subseteq\Pos{A\Qsys}$ by
definition, and that inclusion is all that \emph{First} to \emph{Third}
need. \emph{First}: this is the cone clause of
Theorem~\ref{thm:branch-formula}, and its derivation is short: by
\eqref{eq:fK} and \eqref{eq:traces},
$K_{\hat c}=T_{\hat c}+\tr\Theta_{\hat c}$, both summands non-negative for
$\omega\succeq0$; and $K_{\hat c}\le P_{\hat c}$ is
\xref{foundations/lem:pinchmono} applied to $\tau=\cG_{\hat c}(\omega)$,
that lemma being stated on all of $\Pos{}$ and consuming no trace
hypothesis. \emph{Second}: Lemma~\ref{lem:cone-concavity}\,(b),\,(c),
whose hypothesis on $\{F_v\}$, sector-block-diagonal positive semidefinite
operators, is implied by Definition~\ref{def:admissible-declaration}. That lemma, and not
Corollary~\ref{cor:concavity}, is the source of the concavity here: the
corollary's domain is $\Sleq{}$, which a point of $\Feas(\varepsilon)$
need not meet once the normalisation is relaxed, and
Lemma~\ref{lem:cone-concavity} is stated for exactly this situation, its
part (c) being that substitution. \emph{Third}: with those two facts the
proof of Theorem~\ref{thm:convexity} runs verbatim with
$\Feas(\varepsilon)$ in place of $\Feas$, since it consumes affinity of
$P_{\hat c}$, non-negativity, concavity and continuity of $K_{\hat c}$, and
convexity and compactness of the feasible set, which are (b), and nothing
else. \emph{Fourth}: the objective factors through $X\mapsto X^{p}$, by
way of the functional $T$ of \eqref{eq:Tcone}, and $X\mapsto X^{p}$ is
undefined off the cone, so none of \emph{First} to \emph{Third} survives a
relaxation of positivity.

\emph{(c), Choi form.} Every $J\in\Feas^{J}(\varepsilon)$ has
$J_s\succeq0$ per block, the one condition of the set that is not relaxed,
and $\chi_{\Qsys}(J)=\bigoplus_s
(W_{\tau}\otimes\idon{})J_s(W_{\tau}^{\dagger}\otimes\idon{})$ is
sector-block-diagonal and positive semidefinite; hence
$\chi_{\Qsys}\big(\Feas^{J}(\varepsilon)\big)\subseteq\Pos{A\Qsys}$, the
one property of $\Feas(\varepsilon)$ the state-form argument uses. So the
bounds $0\le K_{\hat c}\le P_{\hat c}$, the cone concavity of
Lemma~\ref{lem:cone-concavity} and Theorem~\ref{thm:convexity} all apply
to $O_{\alpha}$ on $\acSet\times\chi_{\Qsys}(\Feas^{J}(\varepsilon))$; and
$O^{J}_{\alpha}=O_{\alpha}\circ(\idmap\times\chi_{\Qsys})$ with
$\chi_{\Qsys}$ \emph{linear}, so convexity and lower semicontinuity pull
back to the convex preimage $\acSet\times\Feas^{J}(\varepsilon)$, a
pullback along a linear map that needs no injectivity, as in Step~2 of the
proof of Theorem~\ref{thm:firstorder}. Convexity and compactness of
$\Feas^{J}(\varepsilon)$ itself are (b). The transfer uses no bijection.

\emph{(c), the range fact.} On $\Feas(\varepsilon)$ with $\varepsilon>0$,
$\sum_{\hat c}P_{\hat c}(\omega)=\tr\omega$ for an exactly admissible
$\{F_v\}$ (Definition~\ref{def:declaration}): both POVMs being complete,
the sum over $\hat c$ of \eqref{eq:Phimap} is
$\tr[(\sum_xM_x\otimes\sum_vF_v)\,\omega]=\tr\omega$, and the relaxed
normalisation no longer pins that to $1$. The proof of
Proposition~\ref{prop:value-range}\,(i) reads a KL divergence between two
probability vectors, so it is not available on the relaxed set; read at
$\varepsilon=0$, on $\Feas$ itself, it gives $\hlow\ge0$ there. In the
Choi form the range fact does not transfer either: it is read on the exact
set $\Feas^{J}(0)$, where $\Phimap$ is a probability vector.

\emph{(c$'$).} \emph{Choi form.} $R$ is Hermitian because each $J_s$ is
positive semidefinite, and $G=\iota^{\dagger}\tau_A\iota
=W_{\tau}^{\dagger}W_{\tau}$ by \eqref{eq:Wtau}. Then $G\succeq0$ and
$\tr G=\tr[\tau_A\,\iota\iota^{\dagger}]=\tr[\Pi_{\iota}\tau_A]
=\tr\tau_A=1$ by \eqref{eq:iota}, so $G$ is a state. In
\eqref{eq:choi-trace-residual} the first equality is the computation in
the proof of Theorem~\ref{thm:choi}. Hence
$\big|\tr[\chi_{\Qsys}(J)]-1\big|=|\tr[GR]|\le\|R\|\,\tr G=\|R\|
\le d_{A_0}\|R\|_{\max}\le d_{A_0}\varepsilon$, using
$|\tr[Y\varrho]|\le\|Y\|\tr\varrho$ for $\varrho\succeq0$ and the
inequality $\|X\|\le\|X\|_F\le d\,\|X\|_{\max}$ for a $d\times d$
matrix $X$, here on the $d_{A_0}\times d_{A_0}$ matrix $R$. This is the
explicit operator-norm conversion that
\xref{foundations/req:norm-conversion} asks for, with its dimension factor
written into the clause, and it is not
a conversion between the spectral and the entrywise relaxation shapes of
Definition~\ref{def:relaxed-feasible-set}: $R$ is measured in
$\|\cdot\|_{\max}$ on both sides of the derivation and no spectral floor
enters it. By \eqref{eq:choi-trace-residual}, $\tr[\chi_{\Qsys}(J)]=1$
exactly when $\tr[GR]=0$. \emph{State form, trusted-source mode.}
$|\tr\omega-1|=|\tr(\omega_A-\tau_A)|\le d_A\|\omega_A-\tau_A\|_{\max}
\le d_A\varepsilon$, by $\tr\tau_A=1$ and $|\tr X|\le d\|X\|_{\max}$ on
the $d_A\times d_A$ matrix $\omega_A-\tau_A$. \emph{Untrusted-source
mode.} There is no clause before it:
$|\sum_s\tr\omega_s-1|\le\varepsilon$ \emph{is} the relaxed coupling
condition (Definition~\ref{def:untrusted-choi}\,(b)).

\emph{(c$''$).} $|\varsigma|\le d_{A_0}\varepsilon$ is (c$'$), and
$\sum_{\hat c}P^{J}_{\hat c}(J)=1+\varsigma$ is the sum computed in the
proof of the range fact above, read at $\chi_{\Qsys}(J)$. The identity
\eqref{eq:penaltyshift-fw} reads only that
$\sum_{\hat c}P_{\hat c}=1+\varsigma$ with every $P_{\hat c}\ge0$, never
where the displacement came from; so the full-alphabet penalty sum is
$\sfD(\mathbf q\|\bar{\mathbf P})-\log_2(1+\varsigma)$ with
$\bar{\mathbf P}$ a probability vector. Writing the objective as in the
proof of Proposition~\ref{prop:value-range}\,(i),
$O^{J}_{\alpha}=\frac{\alpha}{\alpha-1}
\big[\sfD(\mathbf q\|\bar{\mathbf P})-\log_2(1+\varsigma)\big]
+\sum_{\hat c\in\Chkey}q(\hat c)e_{\hat c}$ with every $e_{\hat c}\ge0$ by
\emph{First} of (c); where some $P^{J}_{\hat c}=0$ with $q(\hat c)>0$ the
objective is $+\infty$ and the bound is trivial. So, provided
$d_{A_0}\varepsilon<1$, \eqref{eq:relaxed-objective-floor} holds. The
state form is the same computation with the factor of (c$'$) for its mode.

\emph{(d).} An infimum over a larger set is no larger, so by (a) the first
two inequalities of \eqref{eq:feas-tol-mono} hold;
$\hlow[\Feas(0)]=\hlow$ because $\Feas(0)=\Feas$; and the last step is
Theorem~\ref{thm:reduction} (Corollary~\ref{cor:reduction-constrained} for
a nonempty $\CList$).

\emph{(d$'$).} At $\varepsilon=0$ the two quantities agree by
Theorem~\ref{thm:choi}. The inclusion
$\Feas^{J}=\Feas^{J}(0)\subseteq\Feas^{J}(\varepsilon)$ of (a) is seen
directly as well: a point of $\Feas^{J}$ has
$\sum_s\tr_{\Qsys_s}J_s=\idon{A_0}$, hence coupling residual
$0\le\varepsilon$; has $\tr\chi_{\Qsys}(J)=1$ by
Theorem~\ref{thm:choi}; is per-block positive semidefinite; and satisfies
every $\bowtie_k$ row exactly, hence in its relaxed form. An infimum over a
larger set is no larger, which gives \eqref{eq:feas-tol-mono-J}, the
equality being Theorem~\ref{thm:choi} and the last step
\eqref{eq:feas-tol-mono}.

\emph{(e).} In the general format the frame is
$\iota=\bigoplus_{\mathsf t}\iota_{\mathsf t}$, which satisfies
\eqref{eq:iota}, and $W_{\tau}=\bigoplus_{\mathsf t}W_{\mathsf t}$
(Lemma~\ref{lem:choi-factor}\,(ii)); so
$G=W_{\tau}^{\dagger}W_{\tau}=\bigoplus_{\mathsf t}G_{\mathsf t}
=\iota^{\dagger}\tau_A\iota$ is a state, the coupling residual of
$J_s\coloneqq\bigoplus_{\mathsf t}J_{(\mathsf t,s)}$ is
$R=\bigoplus_{\mathsf t}R_{\mathsf t}$ on
$A_0=\bigoplus_{\mathsf t}A_{0,\mathsf t}$, and
$\|R\|_{\max}=\max_{\mathsf t}\|R_{\mathsf t}\|_{\max}$. With these
identifications each proof above applies as written. (a) and (b) read each
relaxed line as an affine preimage of a closed convex set, the
\tref{party-local} rows among them, and bound the trace from the coupling
clause, $\sum_{\mathsf t,s}\tr J_{(\mathsf t,s)}\le d_{A_0}(1+\varepsilon)$
with $d_{A_0}=\sum_{\mathsf t}d_{A_{0,\mathsf t}}$. (c) reads per-block
positive semidefiniteness, the linearity of $\chi_{\Qsys}$ and
Theorem~\ref{thm:convexity}, whose inputs in the general format are the
convexity of the feasible set and the concavity of $K_{\hat c}$ on the
joint-sector lists (Definition~\ref{def:feasible},
Lemma~\ref{lem:joint-sectors}). (c$'$) reads $G$, $R$ and \eqref{eq:iota}:
$\tr[\chi_{\Qsys}(J)]=\tr[G(\idon{A_0}+R)]=1+\tr[GR]
=1+\sum_{\mathsf t}\tr[G_{\mathsf t}R_{\mathsf t}]$ and
$|\tr[GR]|\le\|R\|\le d_{A_0}\|R\|_{\max}\le d_{A_0}\varepsilon$. (c$''$)
reads (c$'$) and the statistics, and (d), (d$'$) read (a),
Theorem~\ref{thm:reduction} (Corollary~\ref{cor:reduction-constrained})
and Theorem~\ref{thm:choi}, all stated for the general format.
\end{proof}
